\section{Common residuals and moment prices}\label{mc:section}

Fix $0<p<1$, write $q=1-p$ and $v=pq$, and let $B$ be
Bernoulli with parameter $p$. Conditional on $B=i$, let $U$ have
the law of a centered variable $U_i$, where
\begin{equation}\label{mc:support}
 -O_0\le U_0\le a_0,\qquad -a_1\le U_1\le O_1,
 \qquad a_0+a_1\le\Delta<1.
\end{equation}
Here $O_i>0$ and $a_i\ge0$. Put $s_i=\E U_i^2$,
$m_i=\E U_i^3$ and $s=qs_0+ps_1$. We assume that
$Y=B+U-p$ is nonpositive on $\{B=0\}$ and nonnegative on
$\{B=1\}$. Thus $X=Y/\sqrt{v+s}$ has mean zero and variance one.
Throughout this section
\[
 \rho=\E|X|^3,\qquad R_0=v(1-2v)
\]
refer to this actual law.

Let $(B_i,\widetilde U_i)$ be independent copies of $(B,U)$, and put
$S_n=\sum_{i=1}^n(B_i+\widetilde U_i)$. For $1\le k<n$,
write $N=n-k$, $\pi_j=\binom njp^jq^{n-j}$,
$F_k=\sum_{j=0}^k\pi_j$, and
\[
 H_r=\sum_{j=k+1}^{k+r}\pi_j,\qquad d=1-\Delta.
\]
On an auxiliary probability space take independent pairs
$(U_{0i},U_{1i})$, each with the prescribed marginals, and set
$W_j=\sum_{i\le j}U_{1i}+\sum_{i>j}U_{0i}$.
Conditioning on $\sum B_i=j$ gives the same residual law $W_j$.

\begin{lemma}[One coupled staircase]\label{mc:staircase}
Let $t$ be real and choose any slope satisfying
\[
 \eta\ge\max_{1\le r\le N}\frac{H_r}{rd}.
\]
Then
\begin{equation}\label{mc:stoploss}
 \PP\{S_n\le k+t\}
 \le F_k+\eta\left\{t+\E(W_k-t-\pi_k/\eta)_+\right\}.
\end{equation}
If $h>0$ satisfies
\begin{equation}\label{mc:square-gate}
 \frac{H_r}{\pi_k}\le\frac{2rd}{h}+\frac{r^2d^2}{h^2}
 \quad(1\le r\le N),
\end{equation}
then
\begin{equation}\label{mc:clipped-square}
 \PP\{S_n\le k+t\}
 \le F_{k-1}+\frac{\pi_k}{h^2}\E(t+h-W_k)_+^2.
\end{equation}
In particular, writing $V=\E W_k^2=Ns_0+ks_1$ gives
\begin{equation}\label{mc:full-square}
 \PP\{S_n\le k+t\}
 \le F_k+\frac{\pi_k}{h^2}(V+t^2+2ht).
\end{equation}
When $\pi_{k+1},\ldots,\pi_n$ decrease, the least legal stop-loss
slope is $\pi_{k+1}/d$, and it suffices in
\eqref{mc:square-gate} to check $r=1$.
\end{lemma}
\begin{proof}
The support bounds give $W_{k+r}-W_k\ge-r\Delta$ in this coupling.
After discarding losses at counts below $k$, the CDF minus $F_k$
is at most $\E g(t-W_k)$, where
\[
 g(y)=-\pi_k\mathbf1_{\{y<0\}}
       +\sum_{r=1}^N\pi_{k+r}\mathbf1_{\{y\ge rd\}}.
\]
For every $y$, including each jump point,
$g(y)\le\eta y+\eta(-y-\pi_k/\eta)_+$.
Taking expectations proves \eqref{mc:stoploss}.
Likewise $g(y)+\pi_k\le\pi_k(y+h)_+^2/h^2$:
at the $r$th incoming jump this is precisely
\eqref{mc:square-gate}, and between jumps the right side increases.
This proves \eqref{mc:clipped-square}; replacing the positive part
by the full square gives \eqref{mc:full-square}.
Finally $H_r\le r\pi_{k+1}$ in the decreasing case.
\end{proof}

For $t<0$ and $V\le\overline V$, the choice
$h=-t+\overline V/(-t)$ in \eqref{mc:full-square}, when legal,
gives $F_k-\pi_k t^2/(\overline V+t^2)$. All these inequalities
hold for closed CDFs and allow arbitrary atoms in the conditional laws.

The moment identity, cubic endpoint inequality and shared budget below
come from the framework of Rodziewicz--Ko\l{}odziejek~\cite{RK2026};
we retain their allocation explicitly for the conditional prices.
\begin{lemma}[The actual third-moment budget]\label{mc:budget}
The standardized moment satisfies
\begin{equation}\label{mc:rho-identity}
 \rho(v+s)^{3/2}=R_0+3v(s_0+s_1)-qm_0+pm_1.
\end{equation}
If $s>0$, put $\alpha=ps_1/s$. For any positive cap
$a_0+a_1\le\delta$,
\begin{equation}\label{mc:allocation-floor}
 \rho\ge g_\delta(\alpha):=
 \frac{R_0+\{3[p+\alpha(q-p)]-\delta\}s
 +\{(1-\alpha)/\sqrt q+\alpha/\sqrt p\}^2s^2/\delta}
 {(v+s)^{3/2}}.
\end{equation}
The same allocation gives the count variance
\begin{equation}\label{mc:count-variance}
 V_k=s\left\{\frac{N(1-\alpha)}q+\frac{k\alpha}p\right\}.
\end{equation}
\end{lemma}
\begin{proof}
Expand $|Y|^3=(p-U_0)^3$ and $(q+U_1)^3$ in the two groups to
obtain \eqref{mc:rho-identity}. If $Z$ is centered, $Z\le a$ and
$\E Z^2=z$, then
$\E[(a-Z)(Z+z/a)^2]\ge0$ yields
$\E Z^3\le az-z^2/a$. Apply this to $U_0$ and $-U_1$.
Zero caps force the corresponding centered residual and variance
to vanish; omit their zero contributions. Cauchy--Schwarz gives
\[
 -qm_0+pm_1\ge
 -s\{(1-\alpha)a_0+\alpha a_1\}
 +\frac{(\sqrt q\,s_0+\sqrt p\,s_1)^2}{a_0+a_1}.
\]
Bound the first term below by $-s\delta$ and the denominator
above by $\delta$, then use
$v(s_0+s_1)=s[p+\alpha(q-p)]$ to obtain
\eqref{mc:allocation-floor}. Equation~\eqref{mc:count-variance}
follows from independence.
\end{proof}

If $p<1/2$, this floor increases with $\alpha$ for fixed $\delta$.
For any positive affine cap $\delta(\alpha)$ it is convex in
$\alpha$, since an affine square divided by a positive affine
function is convex. At $s=0$ both residuals vanish almost surely,
and \eqref{mc:rho-identity} gives the Bernoulli moment directly.

The support bounds also imply $s_i\le O_i a_i$ by averaging
$(U_i+\ell_i)(r_i-U_i)\ge0$ over the corresponding support
$[-\ell_i,r_i]$. Suppose an upstream resource bound for this law gives
\[
 g(a_0+a_1)+\lambda\{(1-\alpha)\mathfrak h_-
                         +\alpha\mathfrak h_+\}\le\mathfrak R,
 \qquad \lambda=ns,\quad g>0.
\]
Its coefficients must be supplied by the applicable resource theorem.
Combining it with $a_0+a_1\ge s_0/O_0+s_1/O_1$ yields
\begin{equation}\label{mc:signed-plane}
 \lambda\left\{(1-\alpha)
       \left(\mathfrak h_-+\frac g{nqO_0}\right)
 +\alpha\left(\mathfrak h_++\frac g{npO_1}\right)\right\}
 \le\mathfrak R.
\end{equation}
This necessary plane retains the actual allocation and its common noise.

\begin{lemma}[Classical two-moment comparison]\label{mc:two-moment}
For $a>0$, $z\ge0$, let $T_{a,z}$ have mean zero, variance $z$
and atoms $a,-z/a$, with probabilities
$z/(a^2+z),a^2/(a^2+z)$; put $T_{a,0}=0$.
If $U$ is centered, $U\le a$ and $\E U^2\le z$, then
\begin{equation}\label{mc:pinelis}
 \E(U-c)_+^2\le\E(T_{a,z}-c)_+^2\qquad(c\in\mathbb R).
\end{equation}
This comparison is preserved by independent addition. The reference
expectation increases when either $a$ or $z$ increases.
If $U_1,\ldots,U_m$ are independent, centered, bounded above by
one common $b>0$, and have total variance at most $V$, then
\begin{equation}\label{mc:bentkus}
 \E\left(\sum_i U_i-c\right)_+^2
 \le \E(Z_b(V)-c)_+^2,
 \qquad Z_b(V)=\sum_{i=1}^mT_{b,V/m,i},
\end{equation}
where the reference summands are independent.
\end{lemma}
\begin{proof}
Write $l=-z/a$, allowing $\E U^2\le z$ throughout.
If $c\le l$, then $\E(U-c)_+^2\le\E U^2+c^2\le z+c^2$,
the reference value. If $c\ge a$, both
sides vanish. For $l<c<a$, the pointwise majorant
\[
 (u-c)_+^2\le\left(\frac{a-c}{a-l}\right)^2(u-l)^2
 \quad(u\le a)
\]
has a nonnegative quadratic coefficient and equality at both reference
atoms. Its expectation increases with $\E U^2$, proving
\eqref{mc:pinelis}. This is the squared-positive-part case of
Pinelis~\cite[Lemma~3.2 and Theorem~2.1]{Pinelis2006}.
Apply the same comparison to $T_{a_1,z_1}$ with upper bounds
$a_2\ge a_1$ and $z_2\ge z_1$: this proves both parameter
monotonicities without differentiation. Conditioning on the other
independent summands proves preservation by addition. Zero variance
is the zero law; a centered variable with upper endpoint zero is
also zero and may be omitted. Jensen's inequality includes these
boundaries without assigning positive variance to $T_{0,z}$.
After replacement with common upper endpoint $b$, Bentkus's
variance equalization~\cite[Lemma~4.5]{Bentkus2004}, scaled from
$b=1$, gives \eqref{mc:bentkus}.
\end{proof}

Only test-function expectations are compared in this lemma.
The original conditional distributions need not have finite support.
For $K(V,c)=\E(Z_b(V)-c)_+^2$ one also has
\begin{equation}\label{mc:variance-derivative}
 0\le K_V\le1.
\end{equation}
There is an order proof. Both $K(V,c)$ and
$\E(c-Z_b(V))_+^2$ increase with $V$: for the latter use
$-T_{b,z}\overset d=T_{z/b,z}$ and the two parameter monotonicities.
Their sum is $V+c^2$, hence
\[
 0\le K(V_2,c)-K(V_1,c)\le V_2-V_1\qquad(0\le V_1\le V_2).
\]
The finite reference formula is $C^1$ for $V>0$, including atom
activations, so \eqref{mc:variance-derivative} follows.
Zero variance follows directly by Jensen and continuity.

For any finite reference sum $Z$ dominating the squared positive
parts of $W_k$, define
\begin{equation}\label{mc:Q}
 Q_Z(A)=\inf_{h>0}\frac{\E(Z-A+h)_+^2}{4h}.
\end{equation}
The inequality $x_+\le(x+h)_+^2/(4h)$ gives
$\E(W_k-A)_+\le Q_Z(A)$.
Write
\[
 M_j=\sum_{z_i-A+h>0}w_i(z_i-A)^j.
\]
On a fixed active set the value is $(M_2/h+2M_1+M_0h)/4$.
A positive stationary width satisfies $h^2=M_2/M_0$ and has
price $(\sqrt{M_0M_2}+M_1)/2$; its adjacent activation inequalities
must also hold. The objective is convex in $h$, with continuous
derivative at activations. If $Z\le A$, its infimum is zero.
The analogous clipped-square objective and its interior stationary width are
\[
 M_2/h^2+2M_1/h+M_0,\qquad h=-M_2/M_1\quad(M_1<0),
\]
subject also to the legal staircase gate.

\begin{lemma}[Convex reference prices]\label{mc:convex-prices}
For a fixed reference law $Z$, the price $Q_Z$ is nonincreasing,
convex and $1$-Lipschitz in $A$.
\end{lemma}
\begin{proof}
The map $(A,h)\mapsto\E(Z-A+h)_+^2/h$ is a convex perspective,
so partial minimization proves convexity of $Q_Z$.
At a positive optimizer,
$Q_Z'(A)=-(M_0+M_1/h)/2\in[-1,0]$, by
$M_1^2\le M_0M_2=M_0^2h^2$ and $M_0\le1$.
Convexity extends these bounds to all subgradients.
\end{proof}
